\documentclass[10pt,a4paper]{article}
    \usepackage{amsmath}
    \usepackage{multirow}
    \usepackage{url}
    \usepackage{booktabs}
	\usepackage{amsthm}
    
    \newcommand{\thead}[1]{\textbf{\textsf{#1}}}

    \newcounter{romanlist}
    \setcounter{romanlist}{1}
    \newenvironment{inparenum}{\renewcommand{\item}{(\roman{romanlist})\ \addtocounter{romanlist}{1}}\ignorespaces}{\ignorespacesafterend}

	\theoremstyle{definition}
	\newtheorem{definition}{Definition}
	\newtheorem{proposition}[definition]{Proposition}
	\newtheorem{corollary}[definition]{Corollary}

    \title{\LaTeX\ Assignment 2: Environments}
    \author{L. R. Joshua, 20785216}
    \date{2 April 2019}
    \setlength{\belowcaptionskip}{10pt}
    \begin{document}
    \maketitle
    
    \section{Instructions}
    Once again, type exactly what you see here. Remember to replace my initials
    and US number with your own. Also, please look \emph{very carefully} at how initials
    are to be formatted. Remember all references and citations must be automatic.
    
    \section{Quoting, versifying, and defining}
    Wikipedia has the following to say about limericks, and I \textbf{\emph{quote:}}
    \begin{quote}
        A \textbf{limerick} is a kind of a witty, humorous, or nonsense poem, especially
        one in five-line anapestic or amphibrachic meter with a strict
        rhyme scheme (AABBA), which is sometimes obscene with humorous
        intent.
    \end{quote}
    
    I am particularly fond of naughty limericks. However, lest anybody be offended,
    let's settle for one slightly less dangerous:
    \begin{verse}
        There was a young man from Madras\\
        Who had a magnificent ass.\\
        \quad Not rounded and pink\\
        \quad As you probably think---\\
        It was grey, had long ears, and ate grass!
    \end{verse}
    Here, we use the \texttt{verse} environment. For the indents in lines 3 and 4, google for
    what the \verb$\quad$ and \verb$\qquad$ commands do.
    For using the backslash in typewriter type, consider which verbatim commands you have at your disposal.
    If you are slightly confused about the definition of a limerick, maybe I should
    include the following descriptions.
    \begin{description}
    \item[amphibrachic]{A metrical foot in formal poetry; consists of a long syllable
    between two short ones.}
    \item[anapestic]{(also anap\ae{}stic.)  A metrical foot in formal poetry; consists of two
    unstressed syllables followed by one stressed syllable.}
    \end{description}
    
    \section{A spot of math}
    For all of the following, use the simplest \LaTeX\ code you can. In particular, use the
    \AmS-\LaTeX\ extensions.
    
    \begin{table}[tb]
    \begin{minipage}[t]{0.45\textwidth}
        \caption{An alignment and spanning exercise for tabular data.}
        \label{table:abcd}
        \centering
        \begin{tabular}{||c|c|c|c||}
            \hline\hline
            \multicolumn{4}{||c||}{ABCD} \\
            \hline\hline
            A & B & C & D \\
            \hline
            \multicolumn{2}{||l|}{AB} & \multicolumn{2}{|r||}{DC} \\
            \hline
            A & \multicolumn{2}{|c|}{BC} & D \\
            \hline\hline
        \end{tabular}
    \end{minipage}
    \hfill
    \begin{minipage}[t]{0.45\textwidth}
        \centering
        \caption{Another exercise for aligning and spanning tabular data.}
        \label{table:abc}
        \begin{tabular}{|c|c|}
            \hline
            \multicolumn{2}{c}{ABC} \\
            \hline\hline
            \multirow{3}{*}{AB} & C1 \\
            \cline{2-2} & C2 \\
            \cline{2-2} & C3 \\
            \hline\hline
            \multicolumn{2}{c}{ABC} \\
            \hline
        \end{tabular}
    \end{minipage}
    \end{table}
    
    Consider a real function $f(x)$ such that $f(x)$ and its derivative $f'(x)$ are
    continuous on the interval $[a, b]$. The length $s$ of the graph of $f$ between $x = a$
    and $x = b$ can be found as follows. Consider an infinitesimal part $ds$ of the
    curve $s$. Then, according to Pythagoras's Theorem,
    \begin{align}
        ds^2 & = dx^2 + dy^2, \\
        \frac{ds^2}{dx^2} & = 1 + \frac{dy^2}{dx^2}, \\
        ds & = \sqrt{1 + \left(\frac{dy}{dx}\right)^2}dx, \\
        s & = \int_a^b\sqrt{1+[f'(x)]^2}dx.
    \end{align}
    If a curve is defined parametrically by $x = X(t)$ and $y = Y (t)$, then the
    arc~length between $t = a$ and $t = b$ is
    \begin{align}
        s = \int_a^b{\sqrt{[X'(t)]^2 + [Y'(t)]^2}}dt.\label{equations:length_of_arc}
    \end{align}
    Equation (\ref{equations:length_of_arc}) is a consequence of the distance formula, where instead of
    $\Delta x$ and $\Delta y$, we take the limit. A useful mnemonic is
    \begin{align}
        s = \lim\sum_a^b{\sqrt{\Delta x^2 + \Delta y^2}} = \int_a^b{\sqrt{dx^2 + dy^2}} = 
		\int_a^b{\sqrt{\left(\frac{dx}{dt}\right)^2}+\left(\frac{dy}{dt}\right)^2}dt.
	\end{align}

    \section{Floating environments}
    The first two tables are just here to check if you understand alignment and
    column spanning. To understand their placement, consider the following: A
    single \texttt{table} environment may contain more than one caption. So, Tables \ref{table:abcd}
	and~\ref{table:abc} have been placed in a single \texttt{table} float (with the \texttt{tb} 
	position specifier), but in separate \texttt{minipage} environments. Each \texttt{minipage}
	is exactly $\frac{9w_t}{20}$ wide, where $w_t$ is the default text width;
    the predefined length \verb$\textwidth$ will be helpful in this regard.
    They are pushed apart by a rubber length, and the contents of each are centred 
    with respect to the minipage. The captions are aligned by using the 
    optional \texttt{t} argument for each minipage.
    
    As a rule of thumb, \emph{never use} the \texttt{center} environment for centring inside
    a float, since \texttt{center} introduces extra vertical spacing. That said, inside a
    \texttt{table}, the default spacing between the caption and the rest of the contents
    is not sufficient. Set the length \verb$\belowcaptionskip$ appropriately inside each
    affected environment.
    
    \begin{table}[tb]
        \caption{
            HIV prevalence estimates and the number of people living with HIV,
            2001--2011, \emph{Mid-year population estimates}, P0302, 2011,
            Statistics South~Africa.
            }
            \small
            \begin{tabular}{rcccc}
                \toprule
                \thead{Year} & \multicolumn{3}{c}{\thead{Prevalence}} & \thead{HIV population} \\
                \cmidrule{2-4}
                \phantom{a} & \thead{Women, 15--49} & \thead{Adult, 15--49} &
				\thead{Total population} & \thead{(millions)} \\
                \midrule
                2001 & 17.4 & 16.0 & 9.4  & 4.21 \\
                2002 & 17.7 & 16.2 & 9.6  & 4.37 \\
                2003 & 18.0 & 16.2 & 9.7  & 4.49 \\
                2004 & 18.1 & 16.2 & 9.8  & 4.59 \\
                2005 & 18.3 & 16.2 & 9.9  & 4.69 \\
                2006 & 18.9 & 16.6 & 10.2 & 4.87 \\
                2007 & 18.9 & 16.5 & 10.2 & 4.95 \\
                2008 & 18.9 & 16.4 & 10.3 & 5.02 \\
                2009 & 19.1 & 16.4 & 10.4 & 5.13 \\
                2010 & 19.3 & 16.5 & 10.5 & 5.26 \\
                2011 & 19.4 & 16.6 & 10.6 & 5.38 \\
                \bottomrule
            \end{tabular}
            \label{table:hiv}
        \end{table}

    Table \ref{table:hiv} is in its own \texttt{table} float and also uses the \texttt{booktabs} package\footnote{
        The \texttt{booktabs} documentation is normally available as \url{/usr/share/doc/texlive-doc/latex/booktabs/booktabs.pdf}.
        To typeset URLs, use the \texttt{url} package.
    } to draw well-spaced rules. The tabular data is centred with respect to the \texttt{table},
    as well as inside each column. The fourth column presents an alignment problem,
    however, since the first couple of numbers have only two digits; the \verb$\phantom$
    command will be useful here. Also, the table was set in \verb$\small$ type, and the
    headings are bold and sans-serif.

    \section{New environments}
    Now, also
    \begin{inparenum}
        \item define a environment \texttt{inparenum} that
        \item provides for automatically numbered items
        \item inside a paragraph.
    \end{inparenum}
    For example, the previous sentence was entered as:\footnote{
        The code itself was typeset in a \texttt{verbatim} environment inside a \texttt{quote} environment.
    }
\begin{quote}
\begin{verbatim}
Now, also
\begin{inparenum}
\item define a environment \texttt{inparenum} that
\item provides for automatically numbered items
\item inside a paragraph.
\end{inparenum}
\end{verbatim}
\end{quote}
    To accomplish this, define a new counter and renew the \verb$\item$ command---
    inside the definition of \texttt{inparenum}, otherwise all other items will 
    also be affected---so that this counter is incremented and displayed as a lowercase 
    roman numeral. Suppress any extra spaces that \verb$\newenvironment$ introduces 
    by using the \verb$\ignorespaces$ command in the ``begin'' block and \verb$\ignorespacesafterend$
    in the ``end'' block. \emph{You may not include and use any package that provides inparagraph
    list or enumeration environments}.

	\section{Theorems and so on}
	The following results are from Koblitz \cite[pp. 34--35]{koblitz}.
	\begin{definition}
		A \emph{generator} $g$ of a finite field $\mathbf{F}_q^*$ is an element
		of order $q - 1$; equivalently, the powers of $g$ run through all of the
		elements of $\mathbf{F}_q^*$.
	\end{definition}
	\begin{proposition}
		Every finite field has a generator. If $g$ is a generator of 
		$\mathbf{F}_q^*$, then $g^j$ is also a generator if and only if $gcd(j, q - 1) = 1$.
		In particular, there are a total of $\phi(q - 1)$ different generators
		of $\mathbf{F}_q^*$.
	\end{proposition}
	\begin{corollary}
		For every prime $p$, there exists an integer $g$ such that the powers
		of $g$ exhaust all nonzero residue classes modulo $p$.
	\end{corollary}
	\begin{proposition}
		There exists a sequence of primes $p$ such that the probability
		that a random $g \in \mathbf{F}_q^*$ is a generator approaches zero.
	\end{proposition}

	Note that there are three different theorem-like environments, but that they
	are all numbered with the same counter. When you are typing these results,
	ensure that you use \emph{math mode everywhere necessary}, even though the default
	typeface of the theorem environments is italics.

	Finally, note that Lamport \cite{lamport} defines theorem-like environments to be
	typeset in italics, which makes them a bit difficult to read. Therefore, use
	the	\verb$\theoremstyle$ command from the \texttt{amsthm} package to
	get normal roman text.

	\begin{thebibliography}{0}
		\bibitem{koblitz} Koblitz, Neal. 1994. \emph{A Course in Number Theory
			and Cryptography}. 
			Second edition. New York: Springer.
		\bibitem{lamport} Lamport, Leslie. 1994. \emph{\LaTeX: A document
			preparation system -- user's guide and reference manual}.
			Reading, Mass.: Addison-Wesley Long Inc.
	\end{thebibliography}

    \end{document}
